\scanpage{5}
\par\noindent\begin{minipage}[t]{.1\linewidth}
\figure{\begin{tikzpicture}[x=.34cm,y=.34cm]
\draw[gray!65,line width=.65pt] (0,0) -- (.4,.5) (.65,.85) -- (1.5,2.08) (1.74,2.4) -- (2.58,3.52);
\draw[gray!65,line width=.65pt] (1.65,-.62) -- (1.9,-.27) (2.17,.09) -- (2.82,.91) (3.62,2.02) -- (3.94,2.43);
\draw[gray!35] (2.99,-.46) .. controls (3.86,-1.34) and (4.05,-.25) .. (4.39,1.11);
\end{tikzpicture}}
\end{minipage}\hspace{.012\linewidth}%
\reversemarginpar
\noindent\tikz[baseline=(num.base)]{\node[draw,circle,inner sep=.6pt] (num) {\shortstack{25\\[-.5ex]\scriptsize$\times$}};}\quad Le produit de deux cycles alg. èquiv à $0$ est-il rat. équiv. à $0$ ? (``pb du carré''). Un $0$-cycle \margin{faux, cf Mumford...\enspace\tikz[baseline=-.4ex,x=.12cm,y=.12cm]{\draw[gray] (.5,2.1) .. controls (-.1,.4) and (0,-.3) .. (.65,-.7) (.62,-.9) -- (.65,-2.3);}}% slanted
qui est alb. èquiv. à $0$ est-il rationnellement èquivalent à $0$ ? Question de l'\unsure{équivalence|èquivalence} rationnelle des $0$-cycles sur une surface.
\normalmarginpar

\noindent 26\quad Théorie de Rosenlicht \unsure{relative|reletive} (sur base quelconque)

\noindent 27\quad \unsure{Adèles|Adéles} des schémas de dim, quelconque. Jacobiennes locales,

\noindent\tikz[baseline=(num.base)]{\node[draw,circle,inner sep=1pt] (num) {28};}\quad Construction d'une théorie abstraite des motifs sur prèschémas de type fini sur $\mathbb{Z}$ (construction catégorique générale),

% archive sheet number: 4

